\chapter[Comparison of the Live-Attenuated Yellow Fever Vaccine 17D-204 Strain
	to Its Virulent Parental Strain Asibi By Deep Sequencing]{Comparison of the Live-Attenuated Yellow Fever Vaccine 17D-204 Strain
		to Its Virulent Parental Strain Asibi By Deep Sequencing\footnote{This is a pre-copyedited, author-produced PDF accepted for publication in \emph{ The Journal of Infectious Diseases}. The version of record, (Andrew Beck et al., Comparison of the Live Attenuated Yellow Fever Vaccine 17D-204 Strain to Its Virulent Parental Strain Asibi by Deep Sequencing. The Journal of Infectious Diseases(2014) 209(3):334-344) is available online at http://jid.oxfordjournals.org/content/209/3/334. \parencite{Beck:2014cm}}}

\section{Abstract}
The first comparison of a live \gls{rna} virus vaccine to its \gls{wt} parental strain by deep sequencing is presented, using as a model the \gls{yfv} live vaccine strain 17D-204 and its \gls{wt} parental strain, Asibi. A commercial lot of 17D-204 vaccine and low passage collection isolate of the Asibi strain were obtained and subjected to sequencing analysis by massively parallel methods to compare population variability along multiple scales of the viral genomes. A modeled exploration of small frequency variants was performed to reconstruct plausible regions of mutational plasticity. Overt quasispecies diversity was observed to be a feature of the parental strain, whereas the live vaccine strain by contrast lacks diversity by multiple independent measurements. A lack of putative attenuating mutations in the Asibi population relative to that of 17D-204 was observed, demonstrating that the vaccine strain was derived by discrete mutation of Asibi and not by selection of genomes in the \gls{wt} population. Analyses such as these of attenuated viruses improve our understanding of the molecular basis of vaccine attenuation and provide critical information on the stability of live vaccines and the risk of reversion to virulence, and indicate that relative population structure is a plausible correlate of attenuation for live viral vaccines.
\section{Introduction}
All of the live attenuated vaccines in use today were derived empirically, and our understanding of molecular basis of attenuation is often rudimentary. Deep sequencing offers the opportunity to investigate population structures of live attenuated vaccine strains, contributions of these features to attenuation and stability, and the potential of reversion to virulence \parencite{Neverov:2010ko,Kennedy:2011ff,Luciani:2012em}. The paradigm for \gls{rna} viruses is that they exist as highly diverse ``quasispecies'' populations, and that measurable features of viral diversity are plausible correlates to virulence and pathogenicity \parencite{Vignuzzi:2006ju,Holland:1992cq}. In this context, the \gls{yfv} vaccine was investigated as a model. To date this technology has been applied to vaccines for \gls{dna} viruses, but not to vaccines for \gls{rna} viruses \parencite{Szpara:2011gg}. The live attenuated \gls{yfv} vaccine strain 17D was derived by Max Theiler and coworkers in 1937, and is a critical tool for control of human \gls{yfv} infection \parencite{Theiler:1937wy}. Despite the production of over 550 million doses of vaccine in the last 70 years, our understanding of how the vaccine is attenuated or how the protective immune response is elicited are very limited. Incomplete deep sequencing of the 17D-204 and other live \gls{rna} virus vaccines has been performed at low coverage, without achieving full assembly of the viral open reading frame \parencite{Victoria:2010kf}. The presented results offer the first comprehensive coverage depth ($>$1000 mean reads/nucleotide position aligned), permitting direct comparison of low-frequency variants between the rare, archived parental strain and commercial vaccine derivative. 

The vaccine was derived from the \gls{wt} strain Asibi, which was isolated in 1927 from the blood of a male Ghanaian patient of the same name. This parental virus was empirically attenuated by serial passage 18 times in mouse embryo tissue, 58 times in minced whole chick embryo tissue, and finally 128 times in minced whole chick embryo without nervous tissue \parencite{Lloyd:1936vn,Theiler:1937wy}. This ``17D'' strain was selected from a panel of other candidates for its desirable attenuated properties, which include loss of both viscerotropism and neurotropism in mammalian models \parencite{Monath:2013cu}. The 17D-204 vaccine substrain was derived from the original 17D strain at passage 204, and is now presently manufactured in the United States, France, Senegal, Russia, and China \parencite[278]{WorldHealthOrganization:2010wo}. Standardization of 17D-derived vaccine is achieved by a seed-lot system, ensuring that all production lots are derived from qualified seeds by one passage in embryonated chicken eggs \parencite[251]{WorldHealthOrganization:2010wo}. By consensus, comparison of Asibi and 17D-204 genomes identified 67 nucleotide differences, encoding for 31 amino acid substitutions \parencite{Hahn:1987tw,Dupuy:1989wb}. Reasonably, it is hypothesized that the historical dataset of mutations differentiating Asibi and 17D-204 represent determinants of mammalian pathogenicity \parencite{Lee:2008gd}. However to date, specific amino acids contributing to the attenuation of 17D have not been elucidated. The presented studies utilize a bioinformatics approach to compare nucleotide diversity of 17D-204 and Asibi, under a hypothesis that viral subpopulations distinguish the \gls{wt} and live-attenuated strains. 
\section{Results}
\subsection{Assembly and Inspection of Viral Genomes} 
Nucleotide counts of five randomly downsampled Asibi strain alignments were averaged and floor-rounded along all sites, and the median coverage of this alignment was 5531 (mean=5562.2, \gls{sd}=2054.1). Median coverage depth for the 17D-204 alignment was 5144 (mean=5682.1, \gls{sd}=2434.6). 
\subsection{Consensus Sequences}
Consensus sequences generated by plurality from \emph{de novo} assemblies were compared to reference Genbank sequences of Asibi and 17D-204 vaccines, inspecting the viral polyprotein and \glspl{utr}. Following comparison of our \emph{de novo} assembled Asibi strain consensus to that of the published Asibi sequence (Genbank accession AY640589), recovery of nucleotide positions 11 through 10843 was observed, with the presence of one silent consensus mutation at nucleotide position A2704G. An independently produced \emph{de novo} consensus sequence of 17D-204 was generated from the same commercial lot of vaccine stock as was used in this study and was found to be identical to the consensus sequence obtained here (Genbank accession JX503529.1). Recovery of nucleotide positions 14 through 10850 for the 17D-204 consensus was observed. Relative to a 17D-204 sequence published by Xie \parencite{Xie:1998wi} (Genbank accession AF052437.1), two nonsynonymous mutations were observed at positions U7496C and C7497U, both independently coding for the same amino acid substitution \gls{ns}4B S67P. 
\subsection{Heuristic Identification of Variants} 
For the Asibi strain alignment, 55 variant sites meeting the inclusion criteria were observed, distributed over the entire genome [Table \ref{tab_variants_Asibi} and Appendix Table \ref{Asibi_full}]. Of these, 9 variants encode amino acid substitutions. For the 17D-204 alignment, 32 variant sites meeting the inclusion criteria were observed; of these, 5 encode amino acid substitutions [Table \ref{tab_variants_17D} and Appendix Table \ref{17D_full}]. 
\subsection{Comparison of Quasispecies Diversity} 
Estimates of diversity were compared at multiple scales [Figure \ref{diversity_nucleotides}], For \gls{se}, median values for Asibi and 17D-204 strains were 0.0179 and 0.0170, respectively, and strains were significantly different (U=61791383, p=$1.60 \times 10^{-9}$) [Figure \ref{boxplot_diversity_whole}]. For Simpson’s $1-D$, median values for Asibi and 17D-204 strains were 0.0047 and 0.0044, respectively, and strains were significantly different (U=61095895, p=$7.81 \times 10^{-7}$). Using only the vaccine mutation sites described by Hahn \parencite{Hahn:1987tw}, median entropy values for Asibi and 17D-204 strains were 0.0194 and 0.0214, respectively, and strains were not significantly different (U=1728, p=.127). Median Simpson’s $1-D$ for the mutation sites was 0.005 and 0.006 respectively, and strains were not significantly different (U=2611, p=0.62) [Figure \ref{boxplot_diversity_hahn}]. Significant strain differences were observed for local comparison of 3$^{\prime}$ untranslated regions, and for viral genes \gls{c}, \gls{prm}, \gls{e}, \gls{ns}1, \gls{ns}2A, \gls{ns}2B, \gls{ns}3, \gls{ns}4B, and \gls{ns}5 [Table \ref{Diversity_stats}]. Quantile-quantile comparisons demonstrate inequality of distribution for error rates, indicating a greater relative presence of high outlying variant frequencies for the Asibi strain. This pattern was recovered both along full genomes and individual genes [Figure \ref{QQ_whole},\ref{QQ_genes}].
\subsection{Intersection of Population Structure For Asibi and 17D-204}
Eight of the previously observed mutant positions differentiating Asibi and 17D-204 strains share considerable variant identity between parent and vaccine, but only in the Asibi alignment. Variability at these shared sites is greater for the Asibi strain alignment, while broadly these sites are fixed to a homogenous identity in the 17D-204 alignment [Figure \ref{marginal_columns}]. Of these eight highly overlapping variants identified by heuristic criteria, the mutation A5153G, present in both viruses, codes for the only amino acid substitution, Ile195Val in the \gls{ns}3 protein (1.39\% Val, 98.6\% Ile in 17D-204), with approximate reverse frequency in the Asibi strain population (98.0\% Val, 1.89\% Ile). 
Variant populations recovered from the Asibi strain intersect with 17D-204 to a greater extent by percentage than vice versa [Figure \ref{intersect_freq}]. Considering percent variant frequency at 74 (67 fixed and 7 clonal) sites differentiating Asibi and 17D-204 \parencite{Hahn:1987tw}, Asibi identity intersects with the consensus 17D-204 identity at 2.89 percent of total coverage in the parental read set (n=11196), while conversely 17D-204 identity intersects with the Asibi consensus at 0.42 percent of total coverage in the vaccine read set (n=1595) [Figure \ref{intersect_freq}]. 

\subsection{Recovery of Small Variants By A Paired-Test Model}
A total of 157 mutations were recovered by this criterion, after excluding insertions, deletions, and stop codons. 109 of the mutations were noncoding, of these 53 represent four-fold degenerate sites, while 48 represent amino acid substitutions [Figure \ref{scatter_types}]. Of nonsynonymous substitutions recovered, 2 were observed in \gls{prm}, 17 in \gls{e}, 4 in \gls{ns}1, 6 in \gls{ns}2A, 3 in \gls{ns}2B, 5 in \gls{ns}3, 1 in \gls{ns}4A, 5 in \gls{ns}4B, and 4 in \gls{ns}5. Significantly represented frequency changes (p$<$0.5) were clustered in specific regions of the viral genome [Figure \ref{manhattan_scatter},\ref{manhattan_peaks}]. Mutations recovered from the paired test model were compared both to the consensus (50 percent plurality) sequences of Asibi and 17D-204 generated \emph{de novo}, and also to lists of differentiating mutations observed previously. The paired-test model recovered all of the 63 of 64 mutations that differentiate the \emph{de novo} assembled sequences by consensus, and 15 (11 plus 4 clonal sites) less than the 17D-204 mutation set described by Hahn, or 79.8 percent recovery \parencite{Hahn:1987tw}.

\section{Discussion}
\subsection{The Wild Type Parental Strain Asibi Is A Diverse Quasispecies}
The \gls{wt} Asibi virus was found to consist of diverse quasispecies, as would be expected of a \gls{rna} virus. By contrast, the 17D-204 vaccine strain population was homogeneous; very limited evidence was found for the existence of \gls{wt} nucleotide identity within subpopulations of the vaccine strain. Production of 17D-204 occurs with some variability due to proprietary methods of manufacture, therefore it is important to ascertain the population identity of 17D-204 at the point of use, and later any effects of these mutations on the viral phenotype in experimental systems. Because analysis of viral quasispecies structure is expected to inform predictive models of pathogenicity and attenuation \parencite{Luciani:2012em}, reconstruction of vaccine populations is a matter of public health importance. The standardized and attenuated properties of 17D-204 provide an excellent framework to construct such a comparative model.

Absolute quasispecies diversity was estimated by diversity indices; Asibi and 17D-204 differed by nonparametric comparison, a pattern that was recapitulated at the level of multiple discrete genes. Quantile-quantile comparisons of error rates demonstrated a greater presence in the Asibi strain of variable sites at the high ranges of the total error distribution [Figure \ref{QQ_whole},\ref{QQ_genes}]. As expected, these methods demonstrate a restricted pattern of variability for both genomes, but especially for the 17D-204 vaccine. This admits the possibility that a small number of large entropy values are present in an otherwise predictably distributed set, an expectation frequently utilized in construction of phylogenetic models of sequence divergence that accommodate among-site differences in mutation rate \parencite{Yang:1993wg}. This property of the sequencing dataset is consistent with a hypothesis that site variants with greatest mutation rates are localized to a small proportion of tolerated sites.

\subsection{Variability in Transition to The 17D Population is Localized}	
In a paired-samples test, significant changes in variant frequencies were localized to several regions of the genome hypothesized to influence virulence and pathogenicity [Figure \ref{manhattan_scatter},\ref{manhattan_peaks}]. Clustering of variants by sliding window average reconstructed clusters of mutations that differentiate Asibi and 17D-204 by consensus, of which the highest peaks (indicating higher density of significant mutations) are located at sites in the viral genes \gls{e}, \gls{ns}2A, \gls{ns}3, \gls{ns}4B, and \gls{ns}5. This pattern is rationally expected. Briefly of these, the flavivirus \gls{e} gene encodes the major external structural element of the virus, of which domain III (nucleotides 1851-2168, containing the largest reconstructed peak) contains immunodominant epitopes and influences adhesion to the extracellular matrix \parencite{Lee:2008gd,Sil:1992vsa}. The \gls{ns}2A gene has been implicated as a factor in the assembly of flaviviruses; mutational analysis of \index{Viruses, Miscellaneous!\Acrfull{kunv}}\gls{kunv} showed influence of the protein over particle egress and intracellular membrane ultrastructure development \parencite{Leung:2008cv}. \gls{ns}3 encodes serine protease, helicase and \gls{ntp} triphosphatase activities \parencite{Bera:2007bc,Luo:2008bg}. The \gls{ns}4B gene of flaviviruses (peak at ~7171) has been implicated as an \gls{ifn} antagonist \emph{in vitro} following infection by \index{Viruses, Miscellaneous!\Acrfull{denv}}\gls{denv}, \index{Viruses, Miscellaneous!\Acrfull{wnv}}\gls{wnv}, and \gls{yfv} \parencite{MunozJordan:2005ia}. \gls{ns}5 encodes methyltransferase and \gls{rna}-dependent \gls{rna} polymerase functions \parencite{Zhou:2007bb,Zou:2011kt}.

\subsection{Wild-Type Sequence Content is Not Contained in the Vaccine Population}
It is especially significant that variants in the vaccine overlap the virulent Asibi genotype at low or undetectable frequencies. The \gls{wt} strain contains a subpopulation with representation of vaccine strain sequence identity. The relationship is unidirectional, meaning that the vaccine does not contain low populations of \gls{wt} variants at the coverage depths achieved. It is not possible to assess linkage of variants beyond read pair distance. However, even independent recovery of parental identity variants was not observed in the vaccine alignment, discounting even a low probability that multiple virulence-determinant sites would be linked to a single genome. The apparent fixation of multiple sites in the vaccine strain indicates an influence of selection pressures exerted at discrete, subgenomic scales. This fixation of the population may in part explain the excellent safety record of the vaccine, although mutation of the virus in vaccinees cannot be excluded. However, these data are supportive of the current opinion that serious adverse events giving rise to a \gls{wt} viscerotropic disease phenotype are a consequence of host effects rather than of specific virulent properties of the vaccine \parencite{Barban:2007hi}.  This also supports an interpretation that the genomic stability of YF--Vax\textsuperscript{\textregistered}  at the point of use is stably maintained by the seed-lot system.

\subsection{Variable YFV Sequence Content is Shared by Other Models of Attenuation}
Intersection of the recovered low-frequency \gls{snp} populations with other models of \gls{yfv} attenuation was observed [Figure \ref{venn_parent},\ref{venn_mutant}]. For the Asibi alignment, 12 sites share variant identity with two other models of \gls{yfv} attenuation. The first of these models, the \gls{fvv}, was isolated concurrently with Asibi, and then passaged 128 times in mouse brain to yield the \gls{fnv} vaccine strain, which was distributed in francophone Africa until high rates of post-vaccinal encephalitis in recipients below 14 years of age prompted discontinuation of \gls{fnv} in 1982. Additionally, passage of the Asibi strain in HeLa cells was performed to produce a virus with similar phenotypic characteristics to 17D, namely reduced virulence in both mice and cynomolgous macaques, and loss of vector competence in \emph{Aedes aegypti} \parencite{Barrett:1990te,Dunster:1999ij}.
Sequences of these attenuation models were compared to the reconstructed variants, recovering some similarities [Tables \ref{tab_variants_Asibi},\ref{tab_variants_17D}, Appendix \ref{Asibi_full} \ref{17D_full}]. The Asibi \gls{e} gene variant nucleotide A1054C (\gls{e} protein residue E27H, 9.74\% H) was previously observed to arise in the attenuated HeLa p6 strain of \gls{yfv} \parencite{Dunster:1999ij}, plus two additional silent variants G1000A and A3274G. The Asibi \gls{e} protein variant C2193U  (\gls{e} protein residue A407V, 13.11\% V) was extinguished by consensus during attenuating passage of \gls{fvv}(V) to produce \gls{fnv}(A) \parencite{Wang:1995vu}. Four silent variants intersect with the \gls{fnv} strain. For the 17D-204 alignment, one silent variant C10367U (1.58\%U) was observed in the 3'\gls{utr}, intersecting the consensus identities of both \gls{fnv}(U) and Asibi(U).

\subsection{Attenuation of Asibi Was Likely By Discrete Selection}
Of particular concern is whether the empirically selected vaccine genotype is the result of either discrete selection at specific sites or macroscopic selection of a population of parental virus. Selection pressures governing the virulence of Asibi are largely unknown. Such pressures are likely to be diverse in nature, involving tissue tropisms, transmission dynamics, host immune responses, and others. It is possible that some variants represent tolerable plasticity in the vaccine genotype, comprising a multivariate sequence space that is flexibly occupied by the virus \parencite{Vignuzzi:2006ju}. Little is known about the adaptive context of quasispecies distributions in arthropod-borne viruses, however selection for phenotypes that productively infect both arthropod and vertebrate hosts almost certainly propels this feature. Arthropod-borne viruses encounter population stresses during \gls{wt} transmission cycles, which are relieved by single-tissue, \emph{in vitro} passage at the cost of multihost fitness \parencite{Ebel:2011do,Coffey:2008jq}. The attenuating passage of 17D outside of the normative mosquito-primate host cycle plausibly represents such an example.

The virulence and pathogenicity-determining features of specific viral quasispecies have largely not been addressed; however some evidence has been offered to suggest that both replication-dependent and replication-independent mechanisms exist and both may be parsed \emph{in vivo} \parencite{Lancaster:2011hf}. Since original derivation of the 17D vaccine was performed outside of naturally infected systems, it is reasonable to expect not only divergent adaptation to these tissues, but also drift from genotypes under selection in the natural mosquito-primate transmission cycle \parencite{Coffey:2011ba,Coffey:2008jq}. Significantly, we have recovered common sequence identity occupied by multiple contexts of \gls{yfv} attenuation. The distribution and nature of these effects along the \gls{yfv} genome are significant targets for further investigation, with great potential to inform the study of flaviviral pathogenesis and rational vaccine design.
\clearpage


\begin{sidewaystable}[p]
	\centering
	\caption{Coding Variants observed in the Asibi strain alignment, using a 1.00 percent frequency cutoff heuristic.}
	\begin{ThreePartTable}
	\taburowcolors 1{white .. gray!20}
	% Please add the following required packages to your document preamble:
%booktabs table
\footnotesize
\begin{tabu} to \columnwidth {X[l]X[l]X[l]X[l]X[l]X[l]X[l]X[l]X[l]X[l]X[l]}
\rowfont[1]{\bfseries}
\toprule
Position & Base & Variant Base & Codon & Local \newline Codon & Gene & Reference Amino Acid & Variant Amino Acid & Coverage Depth & Variant Reads & \% \newline Variant \\ \midrule
620      & U    & C            & 168   & 47                   & prM  & Y                    & H                  & 8174           & 129           & 1.59                \\
1054     & A    & C\tnote{a} & 312   & 27                   & E    & Q                    & H                  & 6743           & 656           & 9.74                \\
2193     & C    & U\tnote{b} & 692   & 407                  & E    & A                    & V                  & 2640           & 346           & 13.11               \\
2349     & C    & U            & 744   & 459                  & E    & A                    & V                  & 2985           & 121           & 4.07                \\
2481     & A    & G            & 788   & 10                   & NS1  & K                    & R                  & 3356           & 420           & 12.53               \\
3171     & A    & C            & 1018  & 240                  & NS1  & E                    & A                  & 2440           & 56            & 2.32                \\
4416     & G    & A            & 1433  & 79                   & NS2B & S                    & N                  & 6744           & 195           & 1.86                \\
5153     & A    & G\tnote{c} & 1679  & 195                  & NS3  & I                    & V                  & 6042           & 103           & 1.7                 \\
9716     & G    & A            & 3200  & 694                  & NS5  & D                    & N                  & 3918           & 118           & 3.02                \\ \bottomrule

\end{tabu}
	\begin{tablenotes}
		\item[a]Intersection of sequence with the HeLa/P6 strain.
		\item[b]Intersection of sequence with the French neurotropic vaccine.
		\item[c]Intersection of sequence with the 17D-204 strain.
	\end{tablenotes}
	\end{ThreePartTable}
	\label{tab_variants_Asibi}
\end{sidewaystable}
\clearpage

\begin{sidewaystable}[p]
	\centering
	\caption{Coding variants observed in the 17D-204 strain alignment, using a 1.00 percent frequency cutoff heuristic.}
	\begin{ThreePartTable}
	\taburowcolors 1{white .. gray!20}
	% 17D heuristic variants
\footnotesize
\begin{tabu} to \columnwidth {X[l]X[l]X[l]X[l]X[l]X[l]X[l]X[l]X[l]X[l]X[l]}
	\rowfont[1]{\bfseries}
	\toprule
Position & Base & Variant Base & Codon & Local \newline Codon & Gene & Reference Amino Acid & Variant Amino Acid & Coverage Depth & Variant Reads & \% \newline Variant \\ \midrule
263      & U    & C            & 49    & 49          & C    & F                    & L                  & 4591           & 46            & 1                   \\
3027     & A    & G            & 970   & 192         & NS1  & K                    & R                  & 5476           & 60            & 1.1                 \\
5153     & G    & A\tnote{a} & 1679  & 195         & NS3  & V                    & I                  & 5057           & 72            & 1.42                \\
5765     & U    & C            & 1883  & 399         & NS3  & Y                    & H                  & 3037           & 31            & 1.01                \\
9093     & A    & G            & 2992  & 486         & NS5  & E                    & G                  & 2969           & 40            & 1.35                \\ \bottomrule

	\end{tabu}
	\begin{tablenotes}
		\item[a]Intersection of sequence with the Asibi strain.
	\end{tablenotes}
	\end{ThreePartTable}
	\label{tab_variants_17D}
\end{sidewaystable}
\clearpage

\begin{sidewaystable}[p]
	\centering
	\caption{Diversity index comparison for individual genes between Asibi and 17D-204, Using the Mann-Whitney U Test.}
	\begin{ThreePartTable}
		\taburowcolors [3] 1{white .. gray!20}
		% Comparative diversity statistics
\footnotesize
\begin{tabu}to \columnwidth {X[l]X[l]X[l]X[l]X[l]X[l]X[l]X[l]X[l]X[l]}
\rowfont[1]{\bfseries}
\toprule
		\multicolumn{2}{c}{} & \multicolumn{4}{c}{Shannon's Entropy} & \multicolumn{4}{c}{Simpson's 1-D}   \\ \midrule
		Gene              & Length(b)     & \multicolumn{2}{c}{Median}         &                &        & \multicolumn{2}{c}{Median} &    &   \\  \midrule
		&               & Asibi  & 17D-204 & Mann-Whitney \emph{U} & \emph{p}        & Asibi & 17D-204 & Mann-Whitney \emph{U} & \emph{p}        \\ \midrule
		5$^\prime$UTR             & 119           & 0.020  & 0.019   & 7058           & 9.7E-01  & 0.006 & 0.005   & 6615           & 1.0E+00  \\
		C                 & 363           & 0.014  & 0.016   & 60022          & 3.8E-02* & 0.003 & 0.004   & 59513          & 2.4E-02* \\
		prM               & 492           & 0.014  & 0.016   & 111041         & 2.5E-02* & 0.004 & 0.004   & 109925         & 1.3E-02* \\
		E                 & 1479          & 0.016  & 0.017   & 1039852        & 2.0E-02* & 0.004 & 0.004   & 1031152        & 7.1E-03* \\
		NS1               & 1056          & 0.014  & 0.019   & 436915         & 2.2E-16* & 0.004 & 0.005   & 440043         & 2.2E-16* \\
		NS2A              & 672           & 0.022  & 0.017   & 276584         & 9.4E-13* & 0.006 & 0.004   & 273610         & 1.8E-11* \\
		NS2B              & 390           & 0.018  & 0.014   & 89943          & 1.0E-05* & 0.005 & 0.003   & 88781          & 5.2E-05* \\
		NS3               & 1869          & 0.019  & 0.017   & 1927976        & 3.8E-08* & 0.005 & 0.004   & 1899208        & 3.7E-06* \\
		NS4A              & 447           & 0.020  & 0.020   & 99266          & 8.7E-01  & 0.005 & 0.005   & 98106          & 6.4E-01  \\
		NS4B              & 750           & 0.022  & 0.017   & 335441         & 1.0E-10* & 0.006 & 0.004   & 333309         & 5.4E-10* \\
		NS5               & 2715          & 0.020  & 0.017   & 4038540        & 9.9E-10* & 0.005 & 0.005   & 4011186        & 1.7E-08* \\
		3$^{\prime}$UTR             & 510           & 0.018  & 0.015   & 150138         & 3.3E-05* & 0.005 & 0.004   & 142853.5       & 3.3E-05* \\ \bottomrule

	\end{tabu}
		\begin{tablenotes}
			\item[*]Statistically significant at p$<$0.05.
		\end{tablenotes}
	\end{ThreePartTable}
	\label{Diversity_stats}
\end{sidewaystable}
\clearpage


\begin{figure}[p]
	\centering
	\includegraphics[width=\textwidth]{/Chapter_Asibi_17D_NGS_FIGURE/diversity_genewise_FIGURE/diversity_genewise_STANDALONE.png}
	\caption[Diversity indices for each nucleotide position in Asibi and 17D-204, including both Shannon's entropy and Simpsons $1-D$.]{Diversity indices for each nucleotide position in Asibi and 17D-204, including both Shannon's entropy and Simpsons $1-D$. For both diversity indices used, high diversity values are more frequently observed in the \gls{wt} strain, and high values are observed at a greater number of nucleotide positions.}
\label{diversity_nucleotides}
\end{figure} 
\clearpage

\begin{figure}[p]
	\centering
	\includegraphics[width=3in]{/Chapter_Asibi_17D_NGS_FIGURE/diversity_boxplot_whole_FIGURE/diversity_boxplot_whole_STANDALONE.png}
	\caption[Comparison of diversity indices for yellow fever virus Asibi and 17D-204 strains, at multiple scales.]{Comparison of diversity indices for yellow fever virus Asibi and 17D-204 strains, at multiple scales. All statistical comparisons were performed using the Mann–Whitney U test; 2-tailed P values of $<$.05 were considered statistically significant. Downsampled Asibi read sets were averaged and floor-rounded by counts. Box plots show diversity indices computed for every nucleotide position along the entire genomes for Asibi and 17D-204 strains, and compared. All box plots depict the median, first quartile, third quartile, with whiskers showing +/- 1.5 times the interquartile distance. *p$<$.05. Abbreviation: NS, not significant.}
\label{boxplot_diversity_whole}
\end{figure} 
\clearpage

\begin{figure}[p]
	\centering
	\includegraphics[width=3in]{/Chapter_Asibi_17D_NGS_FIGURE/diversity_boxplot_sites_FIGURE/diversity_boxplot_sites_STANDALONE.png}
	\caption[Comparison of diversity indices for yellow fever virus Asibi and 17D-204 strains, at multiple scales.]{Comparison of diversity indices for yellow fever virus Asibi and 17D-204 strains, at multiple scales. All statistical comparisons were performed using the Mann–Whitney U test; 2-tailed P values of $<$.05 were considered statistically significant. Downsampled Asibi read sets were averaged and floor-rounded by counts. Box plots show comparison of diversity indices for the vaccine mutations sites observed by Hahn et al. \parencite{Hahn:1987tw}. All box plots depict the median, first quartile, third quartile, with whiskers showing +/- 1.5 times the interquartile distance. *p$<$.05. Abbreviation: NS, not significant.}
	\label{boxplot_diversity_hahn}
\end{figure}
\clearpage

\begin{figure}[p]
	\centering
	\includegraphics[width=3in]{/Chapter_Asibi_17D_NGS_FIGURE/QQ_entire_FIGURE/QQ_entire_STANDALONE.png}
	\caption[Quantile-quantile comparisons of error rate along all nucleotide positions recovered for Asibi and 17D-204 strains.]{Quantile-quantile comparisons of error rate along all nucleotide positions recovered for Asibi and 17D-204 strains. Downsampled Asibi read sets were averaged and floor-rounded by counts. Error rates were calculated as frequency of base calls at each site not corresponding to the consensus base call of the alignment. Curve shapes show that a greater number of high outlying errors is present in the Asibi strain alignment, compared with that of 17D-204. This pattern is evident for both the entire genome lengths and at the scale of specific genes. Abbreviation: UTR, untranslated region.}
\label{QQ_whole}
\end{figure} 
\clearpage

\begin{figure}[p]
	\centering
	\includegraphics[width=5in]{/Chapter_Asibi_17D_NGS_FIGURE/QQ_genes_FIGURE/QQ_genes_STANDALONE.png}
	\caption[Quantile-quantile comparisons of error rate along all nucleotide positions recovered for Asibi and 17D-204 strains.]{Quantile-quantile comparisons of error rate along all nucleotide positions recovered for Asibi and 17D-204 strains. Graphic shows each gene segment considered separately. Downsampled Asibi read sets were averaged and floor-rounded by counts. Error rates were calculated as frequency of base calls at each site not corresponding to the consensus base call of the alignment. Curve shapes show that a greater number of high outlying errors is present in the Asibi strain alignment, compared with that of 17D-204. This pattern is evident for both the entire genome lengths and at the scale of specific genes. Abbreviation: UTR, untranslated region.}
\label{QQ_genes}
\end{figure} 
\clearpage


\begin{figure}[p]
	\centering
	\includegraphics[width=\textwidth]{/Chapter_Asibi_17D_NGS_FIGURE/fixation_FIGURE/fixation_STANDALONE.png}
	\caption[Intersection of variant identity between Asibi and 17D-204 strain quasispecies is unidirectional, depicted as a barplot of marginal frequency.]{Intersection of variant identity between Asibi and 17D-204 strain quasispecies is unidirectional. Bar plot shows sites of overt quasispecies identity at sites that are fixed in the 17D-204 strain. By relative frequency, Asibi nucleotide identity intersects that of 17D-204 to a greater extent than the converse. Asibi frequencies are averaged for downsampled read sets.}
\label{marginal_columns}
\end{figure} 
\clearpage

\begin{figure}[p]
	\centering
	\includegraphics[width=3in]{/Chapter_Asibi_17D_NGS_FIGURE/intersect_freq_FIGURE/intersect_freq_STANDALONE.png}
	\caption[Intersection of variant identity between Asibi and 17D-204 strain quasispecies is unidirectional, as box and scatterplot.]{Intersection of variant identity between Asibi and 17D-204 strain quasispecies is unidirectional.  Graphic shows a box plot summary of intersection frequency for Asibi and 17D-204 strains, overlaid with scatterpoints, depicting the frequency of sites bearing the quasispecies identity of the opposing viral read set. By relative frequency, Asibi nucleotide identity intersects that of 17D-204 to a greater extent than the converse. Asibi frequencies are averaged for downsampled read sets.}
\label{intersect_freq}
\end{figure} 
\clearpage

\begin{figure}[p]
	\centering
	\includegraphics[width=2.75in]{/Chapter_Asibi_17D_NGS_FIGURE/scatter_types_FIGURE/scatter_types_STANDALONE.png}
	\caption[Scatterplots showing change in nucleotide frequency for paired samples comparison.]{Scatterplots showing change in nucleotide frequency for paired samples comparison. Plot depicts the set of mutations recovered from a binomial model, using the Asibi strain alignment as a control, and 17D-204 as a test case. Filled polygons represent mutants (152) that were recovered concordantly from all downsampled Asibi strain read sets, at p$<$.05. Mutations are classified, relative to the Asibi (control) strain genome, as coding for amino acid substitutions, as silent, and as silent from four-way degenerate sites. The pattern of significant frequency changes shows fixation in the 17D-204 genome relative to the parental strain Asibi.}
	\label{scatter_types}
\end{figure}
\clearpage

\begin{figure}[p]
	\centering
	\includegraphics[width=\textwidth]{/Chapter_Asibi_17D_NGS_FIGURE/manhattan_scatter_FIGURE/manhattan_scatter_STANDALONE.png}
	\caption[Recovery of local plasticity in the YFV genome using the paired=test model, as Manhattan plot.]{Recovery of local plasticity in the \gls{yfv} genome using the paired=test model. Graphic shows a Manhattan plot depicting all possible variant frequency changes in the modeled transition from \gls{yfv} Asibi strain to 17D-204. Polygons represent significant frequency changes, expressed as $\log(1/p)$, following Bonferroni correction.}
	\label{manhattan_scatter}
\end{figure} 
\clearpage

\begin{figure}[p]
	\centering
	\includegraphics[width=\textwidth]{/Chapter_Asibi_17D_NGS_FIGURE/manhattan_peaks_FIGURE/manhattan_peaks_STANDALONE.png}
	\caption[Recovery of local plasticity in the YFV genome using the paired=test model, as a sliding-window summary.]{Recovery of local plasticity in the \gls{yfv} genome using the paired=test model, as a sliding-window summary. Graphic shows a line plot of a sliding-window mean of the most significant variant recovered from each position in the paired test model. The graph permits visualization of mutationally diverse regions of the YFV genome under the selected model, in comparison to the expected mutation set (black points).}
	\label{manhattan_peaks}
\end{figure} 
\clearpage

\begin{figure}[p]
	\centering
	\includegraphics[width=4in]{/Chapter_Asibi_17D_NGS_FIGURE/venns_parent_FIGURE/venns_parent_STANDALONE.png}
	\caption[Venn diagram showing overlap of significantly changed nucleotide frequencies with prototype YFV strains, showing intersection with parental strains only.]{Venn diagram showing overlap of significantly changed nucleotide frequencies with prototype \gls{yfv} strains, showing intersection with parental strains only. A paired-test model recovered patterns of local plasticity in the \gls{yfv} genome. Graphic shows intersection between recovered variant populations with other models of \gls{yfv} attenuation; this is separately performed for consensus sequences of \emph{parent} (this figure) and derivative strains [Figure \ref{venn_mutant}] for HeLa p6 virus and \gls{fnv} \parencite{Wang:1995vu, Dunster:1999ij}.}
\label{venn_parent}
\end{figure}
\clearpage

\begin{figure}[p]
	\centering
	\includegraphics[width=4in]{/Chapter_Asibi_17D_NGS_FIGURE/venns_mutant_FIGURE/venns_mutant_STANDALONE.png}
	\caption[Venn diagram showing overlap of significantly changed nucleotide frequencies with prototype YFV strains, showing intersection with derived strains only.]{Venn diagram showing overlap of significantly changed nucleotide frequencies with prototype \gls{yfv} parental strains. A paired-test model recovered patterns of local plasticity in the \gls{yfv} genome. Graphic shows intersection between recovered variant populations with other models of \gls{yfv} attenuation; this is separately performed for consensus sequences of parent [Figure \ref{venn_parent}] and \emph{derived} (this figure) strains for HeLa p6 virus and \gls{fnv} \parencite{Wang:1995vu, Dunster:1999ij}.}
\label{venn_mutant}
\end{figure}
\clearpage
